% Originally: written for this chapter; no predecessor week file.
% Generator: Claude Opus 5 (High)
% Added 2026-08-11 on the user's explicit instruction, after they asked whether Jordan measure and
% Lebesgue measure are the same thing and then asked for a short digression section on Lebesgue
% theory. This is a deliberate, user-authorised exception to the "enrich existing sections only"
% scope rule in gemini.md, which is recorded there. Nothing here is proved, nothing here is used
% by any later chapter, and the section says both of those things to the reader.

\section{A digression: a glimpse of the Lebesgue integral}
\label{sec:lebesgue_digression}

This course builds one theory of integration: the Riemann integral over Jordan measurable sets. It
is enough for everything in the chapters that follow, and no result later in this document depends
on a single line of this section. Read it as an answer to a question the previous section provokes
and cannot settle.

The question is what \qt{Lebesgue-null} means, given that this course never defines a Lebesgue
integral, and yet the term appears on the problem sheets (\cref{ex:8.2}\textbf{(b)}) and in the
comparison remarks above. It comes from a second theory of integration, and the sanest way to place
it is to see what that theory fixes. \textbf{Nothing below is proved.} The statements are given so
that the vocabulary has somewhere to attach, and so that a reader meeting measure theory later
recognises what it was built for.

\subsection{Two things the Riemann theory cannot do}

\begin{example}[The measurable sets are too few]
\label{ex:lebesgue_defect_sets}
\Cref{prop:algebra_jordan_measurable} says the Jordan measurable sets are closed under finite
unions, intersections and differences. They are \emph{not} closed under countable unions. Each
singleton $\{q\} \subseteq \mathbb{R}$ is Jordan measurable with measure $0$, and yet
\[\mathbb{Q}\cap[0,1] = \bigcup_{k\ge1}\{q_k\}\]
is not Jordan measurable at all: any finite cover of it by cubes also covers its closure $[0,1]$,
so $\mu_{\mathrm{out}} = 1$ while $\mu_{\mathrm{in}} = 0$. \Cref{ex:non_jordan_measurable_open_set}
is worse still, since the set there is \emph{open} and the volume of an open set ought not to be a
question a theory declines to answer.
\end{example}

\begin{example}[Limits escape the class of integrable functions]
\label{ex:lebesgue_defect_limits}
Enumerate $\mathbb{Q}\cap[0,1] = \{q_1,q_2,\dots\}$ and let $f_k$ be the indicator function of
$\{q_1,\dots,q_k\}$. Each $f_k$ is Riemann integrable, having only finitely many discontinuities,
and $\int_0^1 f_k = 0$. The sequence increases pointwise to the Dirichlet function
\[f(x) = \begin{cases} 1, & x \in \mathbb{Q}, \\ 0, & x \notin \mathbb{Q},\end{cases}\]
which is not Riemann integrable, its upper and lower sums being $1$ and $0$ for every partition.

A bounded increasing sequence of integrable functions therefore need not have an integrable limit.
The only exchange theorem available in this course needs \emph{uniform} convergence, and uniform
convergence is a strong hypothesis that the sequences arising in practice routinely fail.
\end{example}

Both defects have the same shape. The Riemann--Jordan theory is built on finite covers by cubes of
one size, and it is stable under finite operations and unstable under countable ones. Lebesgue's
construction changes exactly that.

\subsection{Lebesgue measure}

\begin{definition}[Lebesgue outer measure]
\label{def:lebesgue_outer_measure}
For $E \subseteq \mathbb{R}^n$ the \newterm{Lebesgue outer measure}
\germanfor{äusseres Lebesgue-Mass} is
\[\lambda^*(E) := \inf\left\{ \sum_{i=1}^\infty \vol(Q_i) \;\middle|\;
E \subseteq \bigcup_{i=1}^\infty Q_i,\ Q_i \text{ cubes} \right\},\]
the infimum running over all \emph{countable} covers of $E$ by cubes of arbitrary sizes.
\end{definition}

This is the notion already in use above: $E$ is Lebesgue null exactly when $\lambda^*(E) = 0$, which
is \cref{rem:lebesgue_vs_jordan_null} restated. The defect of $\lambda^*$ is that it is only
subadditive, not additive, so one restricts it to a good class of sets.

\begin{definition}[Lebesgue measurable set]
\label{def:lebesgue_measurable}
A set $E \subseteq \mathbb{R}^n$ is \newterm{Lebesgue measurable}
\germanfor{Lebesgue-messbar} if it splits every test set additively:
\[\lambda^*(A) = \lambda^*(A\cap E) + \lambda^*(A\setminus E) \qquad
\text{for every } A \subseteq \mathbb{R}^n.\]
On this class one writes $\lambda := \lambda^*$ and calls it \newterm{Lebesgue measure}
\germanfor{Lebesgue-Mass}.
\end{definition}

\begin{theorem}[What the construction delivers, without proof]
\label{thm:lebesgue_measure_properties}
The Lebesgue measurable subsets of $\mathbb{R}^n$ form a \newterm{$\sigma$-algebra}
\germanfor{$\sigma$-Algebra}: a family containing $\emptyset$ and closed under complements and
under \emph{countable} unions. It contains every open set, every closed set and every Lebesgue null
set. On it, $\lambda$ is countably additive: for pairwise disjoint measurable $E_1,E_2,\dots$,
\[\lambda\left(\bigcup_{k\ge1}E_k\right) = \sum_{k\ge1}\lambda(E_k).\]
Moreover every Jordan measurable set $E$ is Lebesgue measurable with $\lambda(E) = \mu(E)$, and the
inclusion is strict.
\end{theorem}

\begin{remark}[The answer to the question, in one line]
\label{rem:jordan_is_not_lebesgue}
Jordan measure and Lebesgue measure are \emph{not} the same, and the difference is not a matter of
convention. Lebesgue measure is a strict extension: it agrees with Jordan measure wherever the
latter is defined, so no volume computed anywhere in this document is put at risk, and it is
defined on many more sets. The single structural upgrade is from an algebra to a
$\sigma$-algebra, that is, from finite to countable stability, and both defects of
\cref{ex:lebesgue_defect_sets,ex:lebesgue_defect_limits} dissolve in it.
\end{remark}

\subsection{The integral: partition the range, not the domain}

The Riemann integral chops the \emph{domain} into small pieces and hopes the function is nearly
constant on each. Lebesgue chops the \emph{range} instead, and asks how large the set is on which
the function takes each value. The standard image: to count a heap of coins, a Riemann integrator
picks them up in the order they lie and adds as they come, while a Lebesgue integrator first sorts
them into denominations and multiplies. The second survives a much messier heap.

Concretely, a \newterm{simple function} is a finite combination $s = \sum_{i=1}^N c_i \chi_{E_i}$
with $c_i \ge 0$ and $E_i$ measurable, where $\chi_E$ is the indicator function of $E$; its
integral is declared to be $\int s \,d\lambda := \sum_i c_i \lambda(E_i)$. For measurable
$f \ge 0$ one then sets
\[\int f \,d\lambda := \sup\left\{ \int s\,d\lambda \;\middle|\; s \text{ simple},\ 0 \le s \le f
\right\},\]
and for general $f$ one splits $f = f^+ - f^-$ into positive and negative parts, calling $f$
\newterm{Lebesgue integrable} if $\int\lvert f\rvert\,d\lambda < \infty$.

The Dirichlet function of \cref{ex:lebesgue_defect_limits} is now trivial: it is the indicator
function of $\mathbb{Q}\cap[0,1]$, a null set, so its integral is $0$. Its Riemann sums never
settle because they are forced to look at intervals; its Lebesgue integral is immediate because the
set carrying the value $1$ is negligible.

\subsection{What the extra work buys}

\begin{theorem}[The convergence theorems, without proof]
\label{thm:lebesgue_convergence_theorems}
Let $f_k$ be measurable functions on a measurable $E \subseteq \mathbb{R}^n$.
\begin{enumerate}[label=\textbf{(\alph*)}]
  \item \textbf{Monotone convergence.} If $0 \le f_1 \le f_2 \le \dots$ pointwise and
    $f_k \to f$ pointwise, then $\int_E f_k\,d\lambda \to \int_E f\,d\lambda$.
  \item \textbf{Fatou's lemma.} If $f_k \ge 0$, then
    $\int_E \liminf_k f_k \,d\lambda \le \liminf_k \int_E f_k\,d\lambda$.
  \item \textbf{Dominated convergence.} If $f_k \to f$ pointwise and there is a single integrable
    $g$ with $\lvert f_k\rvert \le g$ for all $k$, then $f$ is integrable and
    $\int_E f_k\,d\lambda \to \int_E f\,d\lambda$.
\end{enumerate}
\end{theorem}

\begin{remark}[Why these are the point of the whole construction]
\label{rem:why_dominated_convergence_matters}
Compare the hypotheses. Ours is uniform convergence, a statement about $\sup_x\lvert f_k(x)-f(x)
\rvert$, which fails as soon as the functions develop a moving spike. Dominated convergence asks
only for pointwise convergence plus one integrable envelope, and the envelope is usually the easy
part. Monotone convergence asks for nothing but monotonicity.

This is where the cost of staying inside the Riemann theory is actually paid in this document.
Uniform convergence is assumed by hand in \cref{sec:ascoli_arzela}, in the treatment of improper
integrals in \cref{sec:improper_integrals}, and again when differentiating under the integral sign
in \cref{sec:feynmans_trick}, where a domination hypothesis is what the argument morally wants.
Each of those becomes shorter and more general in the Lebesgue setting.
\end{remark}

\subsection{How the two integrals compare}

\begin{theorem}[Lebesgue's criterion, and the comparison, without proof]
\label{thm:lebesgue_criterion_riemann}
Let $Q \subseteq \mathbb{R}^n$ be a closed box and $f : Q \to \mathbb{R}$ bounded.
\begin{enumerate}[label=\textbf{(\alph*)}]
  \item $f$ is Riemann integrable if and only if its set of discontinuities is Lebesgue null.
  \item If $f$ is Riemann integrable, then it is Lebesgue integrable and the two integrals agree.
\end{enumerate}
\end{theorem}

Part \textbf{(a)} is the reason a course with no Lebesgue integral still has to say
\qt{Lebesgue null}. It is the sharp characterisation of Riemann integrability, and no notion built
from finite covers replaces it: the Dirichlet function is discontinuous everywhere, a set of
measure $1$, and it fails; a function with countably many jumps is discontinuous on a null set, and
it succeeds.

\begin{ainote}
It is tempting to conclude that the Lebesgue integral simply subsumes the Riemann integral. For
\emph{proper} integrals on a box, \cref{thm:lebesgue_criterion_riemann}\textbf{(b)} says exactly
that. For \emph{improper} ones it is false in both directions, and the standard witness is
\[\int_0^\infty \frac{\sin x}{x}\,dx = \frac{\pi}{2},\]
which exists as an improper Riemann integral (\cref{def:improper_integral}) but not as a Lebesgue
integral, because $\int_0^\infty \lvert \sin x\rvert/x \,dx = \infty$ and the Lebesgue theory
integrates only absolutely. The improper Riemann integral is a limit of integrals and can exploit
cancellation between the humps; the Lebesgue integral is defined through $\lvert f\rvert$ and
cannot. So neither theory contains the other outright, and \qt{Lebesgue is strictly stronger} is a
half-truth worth not repeating.
\end{ainote}

\begin{remark}[What to take away]
\label{rem:lebesgue_digression_summary}
Three sentences are enough to carry. Jordan measure is Lebesgue measure restricted to a smaller,
finitely-stable class of sets, and where both are defined they agree. The Lebesgue theory exists
because that class is too small to be closed under limits, and its payoff is the convergence
theorems rather than any new volume. And the one place its vocabulary genuinely reaches into this
course is Lebesgue's criterion, which is why \qt{Lebesgue null} had to be defined here at all.
\end{remark}

% Generator: Claude Opus 5 (High)
% First use of the `aside' environment, added to main.tex in the same pass. This is the intended
% shape: a genuinely startling fact that no later argument touches.
\begin{aside}[Some sets have no volume at all, and a ball can be doubled]
\label{aside:vitali_banach_tarski}
\Cref{thm:lebesgue_measure_properties} carefully says which sets $\lambda$ is defined on, and never
claims that is all of them. It cannot: there exist subsets of $\mathbb{R}$ that are not Lebesgue
measurable. Vitali's construction takes the quotient $\mathbb{R}/\mathbb{Q}$ and picks one
representative from each class inside $[0,1]$. The resulting set $V$ has countably many disjoint
rational translates that together fill up an interval, so countable additivity would force
$\lambda(V)$ to be simultaneously $0$ and positive. No value works, and $V$ is unmeasurable.

The construction needs the axiom of choice, and this is not a technicality. It is consistent with
the other axioms of set theory that \emph{every} subset of $\mathbb{R}$ is Lebesgue measurable, so
the pathology is imported by the choice principle rather than discovered in the reals.

Pushed into $\mathbb{R}^3$, the same idea gives the Banach--Tarski paradox: a solid ball can be cut
into five pieces which, moved only by rotations and translations, reassemble into two solid balls
of the original size. Nothing is stretched, and no volume is created, because the five pieces are
unmeasurable and simply do not have volumes to add up. The paradox is a statement about the axiom
of choice, not about geometry.
\end{aside}

\subsection{Exercises}

% Generator: Claude Opus 5 (High)
\begin{aiexercise}[Null sets from countability]
\label{ex:ai_lebesgue_null_countable}
Work directly from \cref{def:lebesgue_outer_measure}.
\begin{enumerate}[label=\textbf{(\alph*)}]
  \item Show that a single point of $\mathbb{R}^n$ is Lebesgue null.
  \item Show that every countable subset of $\mathbb{R}^n$ is Lebesgue null, and deduce that
    $\mathbb{Q}^n$ is.
  \item Show that a Lebesgue null set has empty interior. (Compare
    \cref{ex:lebesgue_null_sets_properties}\textbf{(a)}, which is the same statement contraposed.)
\end{enumerate}
\end{aiexercise}

% Generator: Claude Opus 5 (High)
\begin{aiexercise}[An integral the Riemann theory cannot see]
\label{ex:ai_lebesgue_three_valued}
Define $f : [0,1] \to \mathbb{R}$ by
\[f(x) := \begin{cases}
1, & x \in \mathbb{Q}, \\
2, & x \notin \mathbb{Q} \text{ and } x < \tfrac12, \\
3, & x \notin \mathbb{Q} \text{ and } x \ge \tfrac12.
\end{cases}\]
\begin{enumerate}[label=\textbf{(\alph*)}]
  \item Compute $\int_{[0,1]} f \,d\lambda$.
  \item Determine the set of points at which $f$ is continuous, and conclude from
    \cref{thm:lebesgue_criterion_riemann} that $f$ is not Riemann integrable.
\end{enumerate}
\end{aiexercise}

% Generator: Claude Opus 5 (High)
\begin{aiexercise}[Dominated convergence where uniform convergence fails]
\label{ex:ai_dominated_beats_uniform}
For $k \ge 1$ let $f_k : [0,1] \to \mathbb{R}$, $f_k(x) := x^k$.
\begin{enumerate}[label=\textbf{(\alph*)}]
  \item Determine the pointwise limit $f$ of $(f_k)$ and show that the convergence is not uniform.
  \item Exhibit an integrable dominating function and apply
    \cref{thm:lebesgue_convergence_theorems}\textbf{(c)} to compute
    $\lim_{k\to\infty}\int_{[0,1]} f_k\,d\lambda$.
  \item Confirm the answer by evaluating $\int_0^1 x^k\,dx$ directly.
\end{enumerate}
\end{aiexercise}

% Generator: Claude Opus 5 (High)
\begin{aiexercise}[Improper Riemann but not Lebesgue]
\label{ex:ai_improper_not_lebesgue}
Define $f : [1,\infty) \to \mathbb{R}$ by $f(x) := (-1)^{k+1}/k$ for $x \in [k,k+1)$,
$k \in \mathbb{N}_{\ge1}$.
\begin{enumerate}[label=\textbf{(\alph*)}]
  \item Show that the improper integral $\int_1^\infty f(x)\,dx$ converges, and identify its value.
  \item Show that $\int_1^\infty \lvert f(x)\rvert\,dx = \infty$, so that $f$ is not Lebesgue
    integrable.
  \item Why does this not contradict \cref{thm:lebesgue_criterion_riemann}\textbf{(b)}?
\end{enumerate}
\end{aiexercise}
